\section*{第 \ref{ch:g6:water-states} 章 --- 水的各种状态}

\begin{solution}{exo:g6:water-states:1}
熔化（清晨草叶上的霜消融）；凝固（冰箱里的冰块）；
汽化（衣服晾干、烧开的水）；液化（冷镜子上的水雾）；
升华（干燥霜冻中渐渐变小的雪堆）；
气体直接变为固体
（窗玻璃上的霜花）。
\end{solution}

\begin{solution}{exo:g6:water-states:2}
露水：液化。没有积水的缩小雪人：升华，
即固体在寒冷干燥的空气中直接变为水汽。冰箱里的“烟”：液化——
室内水汽在冷空气中
凝成小雾团（是小液滴，不是气体）。
\end{solution}

\begin{solution}{exo:g6:water-states:3}
仍为 \qty{412}{g}：质量在状态变化中保持不变——密封的
瓶子既没有失去，也没有获得任何质量。
\end{solution}

\begin{solution}{exo:g6:water-states:4}
加盖，使水汽无法带着其质量逸出（记录必须保持封闭）；
擦干，因为瓶外附着的水或霜会带入外来质量，
这些质量属于房间，而不属于本次记录。
\end{solution}

\begin{solution}{exo:g6:water-states:5}
并未消失：\qty{200}{g} 的水以水汽形式
经厨房的空气散走——汽化只是把质量移出锅，
而没有移出世界。
\end{solution}

\begin{solution}{exo:g6:water-states:6}
约 \qty{1.1}{L} 的冰——体积增大了约十分之一。质量
保持不变：进入时为 \qty{1}{kg}，出来时仍为 \qty{1}{kg}。
\end{solution}

\begin{solution}{exo:g6:water-states:7}
困在管道里的水结冰时会膨胀，管道和
玻璃瓶一样会被撑裂。园丁会在秋天放空水龙头，
使里面没有残留的水可以膨胀。
\end{solution}

\begin{solution}{exo:g6:water-states:8}
冰用 $1.1$ 升的空间装下一升的物质：
同样的空间里，冰比水轻——而“同样空间里更轻”
正是物体漂浮的原因。
\end{solution}

\begin{solution}{exo:g6:water-states:9}
海水（\qty{1}{kg}，液态）——汽化——水汽
（\qty{1}{kg}）——液化——云中小水滴（\qty{1}{kg}）
——凝固——雪（\qty{1}{kg}）——熔化——
融雪水（\qty{1}{kg}）——河水入海（\qty{1}{kg}）。
共六个阶段，质量始终不变。
\end{solution}

\begin{solution}{exo:g6:water-states:10}
请放心。漂浮的冰块已经排开了与自身\emph{重力}相等的柠檬水；
融化后，恰好回到这部分水，体积缩回原先增大的十分之一，
正好流入冰块原先在液面下占据的那部分空间。
杯子仍然满满的，
一滴也不会溢出。
\end{solution}

\begin{solution}{exo:g6:water-states:11}
这是气体直接变为固体的途径：房间空气中看不见的水汽
碰到结冰的玻璃，直接凝成冰，没有经过液态。
它不叫“凝固”，因为玻璃上从未出现过液态水；
这些水直接从气态而来。
\end{solution}

\begin{solution}{exo:g6:water-states:12}
池塘在冬季通过水面向冷空气散热，因此表层水先降到 \qty{0}{\celsius}，
冰首先在那里形成，并因比同样空间里的水轻，
而浮在上面，留在水面。冰盖（以及雪）像一条缓慢的通道，
它保护下方的水，使其热量
只能缓慢地向上
散失。深层水在整个冬天都保持在冰点以上，
鱼儿在由这种奇特而会膨胀的水为它们筑成的
“屋顶”之下游动。
\end{solution}

\begin{solution}{pb:g6:water-states:1}
\textbf{1.} 汽化（温暖水面上的蒸发）。
\textbf{2.} 仍为 \qty{1}{kg}——质量在状态变化中保持不变。
\textbf{3.} 液化。
\textbf{4.} 在 \qty{0}{\celsius} 时。
\textbf{5.} 约 \qty{1.1}{L}。
\textbf{6.} 质量为 \qty{1}{kg}；体积回到 \qty{1}{L}。
\textbf{7.} $1000 - 250 = \qty{750}{g}$ 继续留在溪流中。
\textbf{8.} 不是：\qty{250}{g} 的水以水汽形式随空气流动，质量完好；
它们最可能的下一步是液化——先凝成云，
再化为雨。
\textbf{9.} $750 + 250 = \qty{1000}{g}$：整千克的水
都回来了。
\textbf{10.} 两次都看不见：刚离开大海时是水汽，
再次位于瀑布水潭上方时也是水汽，
两次都处于气态。
\textbf{11.} 蓬松说的是\emph{体积}：雪是夹杂空气的冰晶，占据很大的空间。
我们标记的那部分水，\emph{质量}始终恰好是
\qty{1}{kg}——“同样空间里轻”不等于“物质少”：
这两个词分别指质量和体积。
\textbf{12.} 汽化、液化、凝固、熔化——
对于蒸发的那四分之一：汽化、液化（降雨）、
重新汇入液态水。用到四条途径，
升华与气体直接变为固体两条途径在本次游览中未被用到。
\end{solution}
